\section*{الفصل \ref{ch:b1:expression-control} --- التحكم في التعبير المورّثي}

\begin{solution}{exo:b1:expression-control:1}
\emph{lacI} (\omterm{def:b1:genomes:gene}{مورّثة} \omterm{def:b1:expression-control:operon}{الكابح})، والمروّج، والمُشغِّل، و\emph{lacZ}،
و\emph{lacY}، و\emph{lacA}، وموقع CAP. بلا لاكتوز: مطفأ (\omterm{def:b1:expression-control:operon}{الكابح}
مرتبط). ومع اللاكتوز والغلوكوز: مشتغل بالكاد (\omterm{def:b1:expression-control:operon}{الكابح} منفصل، لكن لا
AMP حلقي مع CAP). ومع اللاكتوز بلا غلوكوز: مشتغل تمامًا. وبلا
السكرين: مطفأ.
\end{solution}

\begin{solution}{exo:b1:expression-control:2}
الناتج الذي يعمل منتشرًا جزيء ينتشر ويستطيع أن يعمل على أي نسخة من
هدفه في \omterm{def:b1:cell-unit-of-life:cell}{الخلية}: وهو \omterm{def:b1:expression-control:operon}{الكابح}. والموقع الذي يعمل موضعيًا تسلسل من \omterm{def:b1:nucleic-acids:chain}{DNA}
لا يؤثر إلا في \omterm{def:b1:genomes:gene}{المورّثات} المرتبطة به فيزيائيًا: وهو المُشغِّل.
\end{solution}

\begin{solution}{exo:b1:expression-control:3}
الركود نحو \qty{36}{min} (من \qty{2}{h} إلى \qty{2.6}{h})؛ والتضاعف
\qty{20}{min} على الغلوكوز و\qty{30}{min} على اللاكتوز.
\end{solution}

\begin{solution}{exo:b1:expression-control:4}
\omterm{def:b1:expression-control:enhancer}{المعزِّز}: تسلسل منظِّم، بعيد غالبًا، يرفع استنساخ \omterm{def:b1:genomes:gene}{مورّثة} حين ترتبط
به عوامل. \omterm{def:b1:expression-control:enhancer}{وعامل الاستنساخ}: \omterm{def:b1:proteins:peptide}{بروتين} فيه نطاق ارتباط بجزيء \omterm{def:b1:nucleic-acids:chain}{DNA} ونطاق
تنشيط أو كبت ينظّم الاستنساخ. \omterm{def:b1:genomes:gene}{ومورّثة} الصيانة: \omterm{def:b1:genomes:gene}{مورّثة} يُعبَّر عنها
في كل \omterm{def:b1:cell-unit-of-life:cell}{خلية} بمستوى ثابت.
\end{solution}

\begin{solution}{exo:b1:expression-control:5}
$I^-$: بنيوي. و$O^c$: بنيوي. و$I^s$: غير قابل للتحريض.
و$I^-/I^+$: محرَّض (\omterm{def:b1:expression-control:operon}{الكابح} $I^+$ يعمل منتشرًا). و$I^s/I^+$:
غير قابل للتحريض ($I^s$ سائد). و$I^+\,O^c\,Z^- / I^+\,O^+\,Z^+$: محرَّض
--- \omterm{def:b1:genomes:gene}{فالمورّثة} $Z$ الوحيدة العاملة خلف مُشغِّل سليم.
\end{solution}

\begin{solution}{exo:b1:expression-control:6}
اللاكتوز: السكر (المحرِّض) يرتبط \omterm{def:b1:expression-control:operon}{بالكابح} و\emph{يحرره} من
المُشغِّل --- فتُصنع \omterm{def:b1:enzymes:enzyme}{الإنزيمات} حين تحضر ركيزتها. والتربتوفان: \omterm{def:b1:proteins:aminoacid}{الحمض
الأميني} (\omterm{def:b1:expression-control:operon}{الكابح} المساعد) يرتبط \omterm{def:b1:expression-control:operon}{بالكابح} و\emph{يمكّنه} من الارتباط
--- فتُصنع \omterm{def:b1:enzymes:enzyme}{الإنزيمات} حين يغيب ناتجها. فالمسالك الهدمية تُشغَّل
بمدخلها، والبنائية تُطفأ بمخرجها.
\end{solution}

\begin{solution}{exo:b1:expression-control:7}
على الغلوكوز: نمو طبيعي (فالغلوكوز لا يحتاج إلى CAP). وعلى اللاكتوز
وحده: لا نمو تقريبًا --- فبلا AMP حلقي لا يستطيع CAP تنشيط مروّج
اللاكتوز، فيبقى \omterm{def:b1:enzymes:enzyme}{الإنزيم} عند بضعة في المئة من مستوياته المحرَّضة حتى
مع تحرر \omterm{def:b1:expression-control:operon}{الكابح}. وعليهما معًا: نمو على الغلوكوز ثم توقف؛ ولا طور
ثانيًا. \omterm{def:b1:enzymes:enzyme}{والإنزيم} منخفض في كل حالة.
\end{solution}

\begin{solution}{exo:b1:expression-control:8}
يقتضي التوهين أن يكون \omterm{def:b1:gene-expression:ribosome}{الريبوزوم} مترجمًا الرسولَ والبوليميراز ما يزال
يستنسخه، حتى يشكّل موضع \omterm{def:b1:gene-expression:ribosome}{الريبوزوم} \omterm{def:b1:nucleic-acids:chain}{RNA} خلف البوليميراز. وفي \omterm{def:b1:cell-unit-of-life:prokeuk}{حقيقيات
النواة} يكون الاستنساخ في \omterm{def:b1:eukaryotic-cell:nucleus}{النواة} والترجمة في \omterm{def:b1:eukaryotic-cell:organelle}{العصارة الخلوية}، ولا يُترجم
الرسول حتى يُعالَج ويُصدَّر.
\end{solution}

\begin{solution}{exo:b1:expression-control:9}
تعمل \omterm{def:b1:expression-control:enhancer}{المعزِّزات} من بعيد، بأي اتجاه وعلى أي جانب، فلا يمكن أن تعمل
بأن تُقرأ أو بأن تضع البوليميراز مباشرة؛ وإنما تعمل بربط عوامل
تلامس المروّج عبر عروة من \omterm{def:b1:nucleic-acids:chain}{DNA}، ولا تهم فيها المسافة والاتجاه إلا
قليلًا.
\end{solution}

\begin{solution}{exo:b1:expression-control:10}
$5000\times 4096\times 4 = \num{8.2e7}$ \omterm{def:b1:water-small-molecules:atp}{ATP}: أي \qty{0.8}{\%} من
الميزانية. ومعدل النمو أدنى بمقدار \qty{0.8}{\%}: ففي كل جيل تهبط حصة
الطافر بالمعامل $2^{-0.008} = 0.9945$. ومن $1:1$ إلى
$1:99$: $0.9945^n = 1/99$، أي $n = \ln 99/0.00555 = 830$ جيلًا ---
أي بضعة أسابيع من الزرع المتصل.
\end{solution}

\begin{solution}{exo:b1:expression-control:11}
كل نمط خلوي يحمل مجموعة خاصة من عوامل الاستنساخ، ترتبط \omterm{def:b1:expression-control:enhancer}{بمعزِّزات}
مورّثاته وتستقدم أستيلازات تفتح صبغينها، بينما تكون \omterm{def:b1:genomes:gene}{مورّثات}
الأنماط الأخرى مُمثيلة ومرصوصة في صبغين متغاير لا يبلغه عامل. وتحفظ
العوامل تعبير بعضها بعضًا (شبكة ذات حالات ثابتة)، وتُنسخ وسوم
الصبغين عند التضاعف، فترث البنات العوامل والمناطق المفتوحة
والمغلقة معًا. \omterm{def:b1:genomes:genome}{فالمورثوم} واحد؛ أما \omterm{def:b1:genomes:genome}{المورثوم} المتاح فليس واحدًا،
وعدم الإتاحة موروث.
\end{solution}

\begin{solution}{exo:b1:expression-control:12}
تحكم البكتيريا مضبوط على الوسط: فسكر أو \omterm{def:b1:proteins:aminoacid}{حمض أميني} أو حرارة يرتبط
\omterm{def:b1:expression-control:operon}{بكابح} أو \omterm{prop:b1:expression-control:trp}{بعامل سيغما}، وتتم الاستجابة في دقائق وتنعكس بمجرد زوال
الإشارة --- أي إن التعبير يتعقب البيئة. \omterm{def:b1:cell-unit-of-life:cell}{وخلية} الجسم تستجيب هي أيضًا
لإشارات، لكن قراراتها الكبرى اتُّخذت في أثناء النمو وأُقفلت في
الصبغين: فما تستطيع التعبير عنه حددته العوامل التي ورثتها والوسوم
على \omterm{def:b1:nucleic-acids:chain}{DNA} فيها، ولا يتغير إلا بالانقسام وهو غير قابل للانعكاس غالبًا.
فالبكتيريا قارئة للحاضر؛ \omterm{def:b1:cell-unit-of-life:cell}{والخلية} المتمايزة ذات ذاكرة هي هويتها.
\end{solution}

\begin{solution}{pb:b1:expression-control:1}
\textbf{1.} $0.5\times 10^{-3}\times 180 = \qty{0.09}{g/L} =
\qty{9e-5}{g/mL}$.
\textbf{2.} $9\times 10^{-5}/1.5\times 10^{-12} = \num{6e7}$ انقسام.
\textbf{3.} $10^7 + \num{6e7} = \num{7e7}$ \omterm{def:b1:cell-unit-of-life:cell}{خلية} لكل mL.
\textbf{4.} $7 = 2^n$: أي $n = 2.8$ تضاعفًا، و\qty{56}{min}.
\textbf{5.} مطفأ: فاللاكتوز (الألولاكتوز) حرر \omterm{def:b1:expression-control:operon}{الكابح}، لكن الغلوكوز
يبقي AMP الحلقي منخفضًا، فلا يرتبط CAP ويكاد المروّج يسكت؛ ثم إن
النافذة نادرة، فلا يدخل من اللاكتوز إلا قليل.
\textbf{6.} يرتفع AMP الحلقي ويرتبط \omterm{def:b1:proteins:peptide}{ببروتين} CAP، الذي يرتبط بالمروّج.
\textbf{7.} رُفع التحكم السالب لكن التحكم الموجب كان غائبًا: فمروّج
ضعيف بلا CAP نادرًا ما يبدأ (\omterm{def:b1:expression-control:cap}{كبت الهدم}).
\textbf{8.} $5000/1800 = 2.8$ رباعي في الثانية؛ و$2.8\times 4096 =
\num{11400}$ بقية في الثانية.
\textbf{9.} السعة $20\,000\times 20 = \num{4e5}$ بقية في الثانية: أي
\qty{3}{\%} من \omterm{def:b1:gene-expression:ribosome}{الريبوزومات}.
\textbf{10.} $5000\times 4096\times 110\times 1.66\times 10^{-24} =
\qty{3.7e-15}{g}$: أي \qty{2.5}{\%} من \omterm{def:b1:proteins:peptide}{بروتين} \omterm{def:b1:cell-unit-of-life:cell}{الخلية}.
\textbf{11.} $5000/5 = 1000$.
\textbf{12.} لا يستطيع اللاكتوز عبور الغشاء بلا نافذة؛ فالطافر
\emph{lacY}$^-$ عنده \omterm{def:b1:enzymes:enzyme}{الإنزيم} ولا ركيزة له في الداخل، ولا ينمو على
اللاكتوز (ولا يتحرض جيدًا، لأن المحرِّض يُصنع من اللاكتوز داخل
\omterm{def:b1:cell-unit-of-life:cell}{الخلية}).
\textbf{13.} $1.0\times 10^{-3}\times 342 = \qty{0.342}{g/L}$، تساوي
$2\times 180 = \qty{360}{g}$ من الغلوكوز لكل مول، أي \qty{0.36}{g/L}
$= \qty{3.6e-4}{g/mL}$.
\textbf{14.} $3.6\times 10^{-4}/1.5\times 10^{-12} = \num{2.4e8}$
انقسام؛ والنهائي $\num{7e7} + \num{2.4e8} = \num{3.1e8}$ لكل mL.
\textbf{15.} $3.1/0.7 = 4.4 = 2^n$: أي $n = 2.1$ تضاعفًا، و\qty{64}{min}.
\textbf{16.} الغلوكوز \qty{56}{min}، والركود \qty{30}{min}، واللاكتوز
\qty{64}{min}: أي \qty{150}{min} في المجمل.
\textbf{17.} يجب حلمأة اللاكتوز أولًا، وهي خطوة إنزيمية زائدة تحد
سرعتُها إمداد الغلوكوز؛ ثم إن النافذة تستعمل تدرج البروتونات
لإدخاله، وهي كلفة لا يدفعها نقل الغلوكوز.
\textbf{18.} $I^-$: \omterm{def:b1:enzymes:enzyme}{الإنزيم} عند \num{5000} طوال الوقت (بنيوي، متى
نفد الغلوكوز؛ وعلى الغلوكوز يظل CAP يحدّه، فيكون المستوى وسطًا)؛
وركوده أقصر أو معدوم، فيُظهر المنحني توقفًا أصغر؛ وتبطئه الكلفة
قليلًا على الغلوكوز.
\textbf{19.} CAP$^-$: ينمو على الغلوكوز إلى $\num{7e7}$ ثم يتوقف:
\omterm{def:b1:expression-control:operon}{فأوبرون اللاكتوز} لا يمكن تنشيطه، ويبقى \omterm{def:b1:enzymes:enzyme}{الإنزيم} قرب 5 جزيئات، ولا
يُستعمل اللاكتوز أبدًا.
\textbf{20.} على اللاكتوز وحده: لا فرق في مستوى \omterm{def:b1:enzymes:enzyme}{الإنزيم} متى تحرّض
(كلاهما مشتغل تمامًا) ولا فرق يذكر في النمو؛ وإنما ينقص $O^c$ تأخيرَ
التحريض. وعلى الغلوكوز وحده: يصنع $O^c$ \omterm{def:b1:enzymes:enzyme}{الإنزيم} بلا داع (بقدر ما
يسمح CAP)، والنمط البري لا يصنعه.
\textbf{21.} $5000\times 4096\times 4 = \num{8.2e7}$ \omterm{def:b1:water-small-molecules:atp}{ATP}: أي \qty{0.8}{\%}
من $10^{10}$.
\textbf{22.} زمن الجيل أطول بمقدار \qty{0.8}{\%}: فمعدل النمو أدنى
بمقدار \qty{0.8}{\%}، أي $\delta = 0.008$ تضاعف في الجيل.
\textbf{23.} $r = 2^{-0.008} = 0.99447$؛ و$r^n = 0.111$: أي $n =
\ln 0.111/\ln 0.99447 = 2.20/0.00555 = 400$ جيل.
\textbf{24.} على اللاكتوز يكون \omterm{def:b1:enzymes:enzyme}{الإنزيم} مطلوبًا على كل حال، فلا يدفع
الطافر شيئًا ويكسب الركود: فلا يُنتخب ضده بل يُنتخب له أحيانًا؛ ثم
إن السلالات المخبرية تُربى للراحة لا للتنافس.
\textbf{25.} معامل التحريض 1000 (من 5 إلى \num{5000} جزيء)؛ والتخليق
بلا داع يكلف \qty{0.8}{\%} من ميزانية \omterm{def:b1:water-small-molecules:atp}{ATP}؛ وتدوم التجربة نحو ساعتين
ونصف.
\end{solution}
